\begin{center}
\textsc{Gauss} an \textsc{Bessel}.\quad Göttingen, 20.~November 1824.
\end{center}

\par\noindent\hspace*{\parindent}.\ .\ .\ .\ .\ .\quad Ich habe in diesem Jahr 12 Stationen besucht und bin mit den Dreiecken bis an die Weser gekommen (der entfernteste Punkt liegt auf der Garlster Heide zwischen Osterholz und Vegesack).
Die dritten Eckpunkte ließ ich anfangs immer zurück, und als ich nach Mitte August Bremen verließ, hatte ich noch 6 Orte zu besuchen.
An allen ging es noch erträglich, aber am letzten, dem Wilseder Berg, quälte mich das schlechte Wetter so, dass ich trotz eines dreiwöchigen Aufenthalts nicht ganz zu meiner Zufriedenheit fertig wurde. Ich musste schließlich einen Endtermin festsetzen und kam Ende Oktober nach Göttingen zurück.

\par\noindent\hspace*{\parindent}Ich behalte mir vor, Ihnen künftig ausführlicher über das Ergebnis der Messungen, die großen Schwierigkeiten, mit denen ich zu kämpfen hatte (eine davon, das ständige Moorbrennen, wird Ihnen selbst noch im Gedächtnis sein), und einige interessante Erscheinungen zu schreiben, die sich dabei gezeigt haben.\quad.\ .\ .\ .\ .\ .

